\section{S3 compatibility and replicated storage}\label{sec:storage}
Part~5 extends the object-storage product in R3. The native storage API, durable backend and shared fine-grained authorization already ship in R1. Payload encryption in both modes, decision-bound key release and opaque envelope custody are part of R1; R2 adds the MCP service without replacing that foundation. S3 is an additional protocol facade over the same content service. A custom replicated blob backend is a separate, optional replacement for the R1 single-node store. Neither S3 compatibility nor owning a replication engine defines whether the product is object storage. Each extension is admitted independently through the existing backend and content contracts; adding S3 does not require changing the byte-storage engine, and changing the backend does not change client policy semantics.

\subsection{S3 compatibility profile v1}\label{sec:s3}
Implement protocols from public specifications with an explicit small surface. Version~1 supports header-signed SigV4 over TLS, bounded known-length single-chunk PUT, GET and HEAD, logical deletion, explicit historical version operations and ListObjectsV2 with filtered visibility. For a server-mode bucket, PUT bodies are plaintext that the storage role encrypts, and GET returns plaintext after the ordinary checks and key release, so a stock S3 client can read and write it. For a client-mode bucket, PUT bodies are client-produced encrypted packages and ordinary GET returns ciphertext, not transparently decrypted file bytes; a qualified client SDK additionally obtains the authorized signed descriptors and recipient envelope through the content/envelope API and decrypts locally, and a stock S3 client can transfer that ciphertext but does not thereby implement the encrypted-object profile. Bucket creation initially uses the control API. Multipart, server-side copy, presigned requests, browser POST, ACLs, public bucket policies, SigV2, SigV4a, streaming-signature modes and directory-bucket extensions are disabled until separately qualified.
\req{S3P-001}{R3}{T-07}{The server advertises ``S3 API subset, content profile v1''. Unsupported requests receive documented protocol errors, the restrictions appear in client configuration examples and tests, and the platform promises neither default S3 client compatibility nor transparent plaintext reads of client-mode content.}
\req{S3P-002}{R3}{T-07}{The profile requires an explicit SHA-256 payload hash (including the empty-body hash) and rejects \code{UNSIGNED-PAYLOAD}, every \code{STREAMING-*} mode and unimplemented checksum or trailer encodings. For a client-mode bucket, the ciphertext body digest and uploader descriptor binding are verified before publication; SigV4 authenticates the request but does not prove payload AEAD validity or plaintext correctness, and the client verifies AEAD tags after download. For a server-mode bucket, the signed payload hash covers the plaintext body, which the storage role encrypts under \rr{SSE-005}.}
\req{S3P-003}{R3}{T-07}{The SigV4 verifier is a new security-sensitive module built from the documented canonicalization and verification requirements; the AWS Go SDK signer is a client helper, not a server verifier. It validates credential scope, service and region, timestamp and skew, signed-header coverage, request canonicalization, signature (with constant-time comparison), credential state and payload integrity. S3 object paths are not normalized, because duplicated slashes and encoded bytes can distinguish keys.}
\req{S3P-004}{R3}{T-07}{S3 routing uses a dedicated raw handler that preserves exact signing and object-key semantics. Go's \code{ServeMux} cleans and redirects some paths and is not used for S3 routing, although management APIs may use it.}
\reqn{S3P-005}{R3}{T-07}{The compatibility matrix names each tested client (for example the AWS SDK for Go v2, boto3 and rclone) with the exact settings it needs, including request-checksum behavior, because recent AWS SDK releases enable default request checksums and chunked trailers that the profile rejects. Failures under default settings are documented.}
\req{S3P-006}{R3}{T-07}{Protocol tests cover duplicate headers and query parameters, percent encoding, literal plus, Unicode, host ambiguity, missing signed security headers, truncation, oversized bodies and time-bound credentials, and use independently authored cases and external client tests.}
\req{S3P-007}{R3}{T-07}{A read authenticates the request, resolves a committed logical object and version, loads trusted attributes, evaluates read or stat, and opens exactly that immutable ciphertext blob, binding resource generation, metadata revision, mode, profile and crypto version to the decision. A server-mode read decrypts only after the key release of \rr{SSE-003}. Client-mode recipient-envelope retrieval remains separately authorized under Table~\ref{tab:keyops}; an S3 credential is not a decryption key.}
\req{S3P-008}{R3}{T-07}{Conditional writes (\code{If-Match}, \code{If-None-Match}) map to the publication compare-and-swap of \rr{PUB-007}, and the platform documents its precise error mapping to S3 behavior.}
\req{S3P-009}{R3}{T-07}{Multipart and copy enter the compatibility matrix only after dedicated tests. Multipart needs an upload binding, creator and delegation, protected intended metadata, part checks, resource limits and a fresh completion decision that uses the same publication compare-and-swap and rejects revoked authority; retries use intent idempotency. Copy checks the exact source version, source read including historical access, destination create or new-version, and explicit transfer restrictions. Because ciphertext and envelopes bind immutable identities, a logical copy to a new identity cannot simply relabel existing ciphertext: a qualified client must construct the newly bound encrypted package, or the operation remains unsupported.}

\subsection{Replicated blob backend}\label{sec:replication}
When the custom backend is selected, its production topology for object bytes is three full replicas in distinct failure domains with two durable receipts required before publication. The system survives the loss of one replica or domain under the stated placement and durability assumptions and repairs back to three. Erasure coding is deferred.
\req{REP-001}{R3}{T-15}{The writer assigns an immutable blob ID and placement generation and sends identical ciphertext to three selected replicas. Each replica verifies length and digest, flushes bytes and directory state, and returns an authenticated receipt binding tenant, node identity and incarnation, failure domain, placement generation, blob and version ID, digest and length. Distinct approved failure domains are counted, not merely processes on one host.}
\req{REP-002}{R3}{T-15, T-17}{Publication requires two valid receipts and stores their placement bindings; it never publishes after only a local memory or disk acknowledgement. A background repair worker copies the identical ciphertext to complete the third replica and does not independently re-encrypt. Under-replicated objects are visible in health metrics and are repaired before planned removal. Receipts from obsolete placement generations or node incarnations cannot satisfy a new publication, and a failed catalog compare-and-swap leaves an orphan eligible for reconciliation.}
\req{REP-003}{R3}{T-15}{Read admission still uses the authoritative catalog and policy epoch. A read may use one recorded healthy replica of the exact immutable blob, checks expected identity and recorded ciphertext integrity data, and fails over to another replica on corruption or failure. Only the recipient client verifies payload AEAD tags; replicas and storage frontends have no payload decryption keys. A full-object digest is confirmed only after the complete stream. Cached blob availability does not authorize a read when authoritative permission state is unavailable.}
\req{REP-004}{R3}{T-17}{Placement uses stable generated IDs, capacity-aware deterministic selection and failure-domain constraints. The chosen node set and placement generation are persisted and are not recomputed from current cluster membership. Rebalancing writes and verifies a replacement replica before updating placement and retiring the old one.}
\req{REP-005}{R3}{T-17}{Repair state is durable. In-memory queues are bounded wakeups and never the sole recovery record. Work is deduplicated by kind, tenant, blob version and target placement generation, and progress, retry time and terminal or unknown status are persisted. A paced scanner reconciles catalog placements, pending intents, tombstones and recorded replica availability and checks ciphertext integrity against stored digests without decryption keys, bounded by bytes, objects, time and tenant quota. Lost wakeups, restarts and interrupted scans do not erase required repair.}
\req{REP-006}{R3}{T-17}{Adding capacity affects eligible future placements and does not automatically move existing blobs. Planned evacuation is an explicit restartable administrative operation: stop new assignment; inventory referenced placements; copy and verify replacements; compare-and-swap the placement generation; verify the required replica and domain count; then retire old copies. Cancellation stops new work while preserving valid committed placements.}
\req{REP-007}{R3}{T-20}{Maintenance is a registered workload, not an identity-less bypass. An internal plan binds deployment, workload action, tenant and blob, permitted source and target nodes, node incarnation, job and fencing token, and placement generation. Repair may copy validated ciphertext for a recorded placement; it may not publish logical content, obtain plaintext keys or delete unrelated blobs. Bootstrap runs in a separate restricted state, never as an omitted-principal shortcut in ordinary traffic.}
\req{REP-008}{R3}{T-17, T-20}{Physical removal requires a current catalog-authorized retirement or tombstone plan and rechecks live references, applicable retention or hold state and current placement at execution. The same destructive-operation guard serves garbage collection, scans and administrative jobs, and job authority and stale fences are checked at the node boundary as well as in the scheduler. The user who requested maintenance is recorded separately from the executing workload.}
\req{REP-009}{R3}{T-16}{The write quorum is not lowered automatically. A partition that cannot obtain two durable domain receipts or a safe catalog commit rejects writes.}
\req{REP-010}{R3}{T-20}{Production internode RPC authenticates workload and node identity, validates typed scope, placement and fencing, bounds streams and messages, and reuses the same receipt and error types as the local backend. It does not duplicate user-facing policy combination in a second storage engine.}
\req{REP-011}{R3}{T-18}{Encrypted, access-separated backups of immutable objects are part of the restore test of \rr{AVL-005}, and restore reconciles catalog references with blobs to a mutually consistent point or quarantines unmatched entries.}
